\section{Interaction Contract}
\label{sec:contract}

Software contracts commonly guard technical boundaries.  In an intelligent
interface, however, the same decisions also determine what the system may show,
which agent may respond, and what explanation or correction point the user
receives.  We illustrate this connection through one running scenario: a
visitor selects \texttt{brand\_panel3} while the Lounge Host is active.  The
system must preserve the selected panel, determine whether the Lounge Host or
the Reception Agent is responsible, enforce any permitted handoff, invoke the
corresponding capability, and present an answer only while that target--agent
relation remains intact.

\subsection{Interaction-design rationale and failure model}

The contract is not itself a user-interface widget.  It is a typed state and
evidence layer from which the interface derives four forms of transparency:
\emph{reference transparency} (what object the system believes was selected),
\emph{responsibility transparency} (which agent may act and why),
\emph{capability transparency} (what operation is available), and
\emph{outcome transparency} (whether the returned response still matches the
pending interaction).  These goals reflect established human--AI interaction
principles: make system status and capability boundaries visible, communicate
why the system acted as it did, and leave users a clear recovery path when the
system cannot proceed \cite{amershi2019guidelines,yang2020reexamining,weisz2024design}.
They also correspond to explanation questions users may reasonably ask, such as
``What did the system select?'', ``Why did this agent answer?'', ``Why not the
agent I approached?'', and ``Why was the request blocked?''
\cite{liao2020questioning}.

The prototype uses progressive disclosure rather than exposing the complete
contract trace to visitors.  A spatial highlight communicates the current
target in place; a handoff cue names a change of responsible agent; and a short,
neutral guidance message communicates ambiguity, unreachability, or an
unavailable capability.  Detailed identifiers and violated relations remain
available for developer diagnosis.  These mappings are design choices, not
claims that a badge or visual cue is universally preferable to text, speech, a
persistent responsibility panel, or an explicit clarification dialogue.  The
formal contract is presentation-independent and can support any of these
alternatives.  Moreover, explanation should not be treated as an automatic
trust intervention: explanations can increase acceptance without necessarily
improving appropriate reliance \cite{bansal2021whole}.  We therefore frame the
signals as support for transparency, comprehension, and user control, and
evaluate the implemented mapping empirically rather than assuming its benefit.

A fluent response can otherwise mask critical interaction breakdowns, such as
referencing an unintended object, presenting an agent as responsible when the
model does not declare that responsibility, or silently failing a handoff.
Table~\ref{tab:failure-model} maps each cross-layer failure to the contract
boundary and to the corresponding transparency or recovery signal.  The final
column describes the current prototype mapping; another interface may render
the same contract state differently.

\begin{table*}[t]
  \caption{Cross-layer failures, contract decisions, and their interaction-design implications.}
  \label{tab:failure-model}
  \small
  \begin{tabularx}{\textwidth}{@{}p{22mm}Y p{34mm}Y@{}}
    \toprule
    Failure & Cross-layer condition & Contract decision & User-facing transparency and control \\
    \midrule
    Target identity
      & A collider reports an unknown entity, an obsolete model hash, or a
        source identifier shared by several objects.
      & Reject before routing or response generation.
      & Clear or invalidate the selection cue and state that the target could
        not be grounded, rather than answering about a guessed object. \\
    Target ambiguity
      & Several raycast candidates overlap, or a deictic request has no unique
        referent.
      & Reject candidate selection instead of choosing silently.
      & Keep the unresolved state visible and request a new or disambiguated
        selection. \\
    Responsibility drift
      & The active agent is not the unique owner and has no declared direct
        handoff to the owner.
      & Apply the responsibility hierarchy and permit at most one direct
        handoff.
      & Name a valid handoff and its destination; otherwise give a scoped
        refusal that distinguishes ``not responsible'' from a generic failure. \\
    Action shortcut
      & A scenario or client attempts to invoke an endpoint without a matching
        provider-specific capability and runtime binding.
      & Require one exact scenario-step--capability-use--capability--binding--action
        path.
      & Do not call the response model; communicate that the requested
        capability is unavailable for this target and provider. \\
    Evidence mismatch
      & A response names a different target, provider, binding, or action from
        the pending request.
      & Reject presentation and retain the violated relation.
      & Block the mismatching answer while preserving the visible pending
        target, so transport success is not presented as interaction success. \\
    Undeclared handoff
      & The response model proposes an agent that is not an allowed handoff
        target.
      & Validate the proposal before commit and roll back provisional state.
      & Abort the handoff transition and report that the proposed delegation is
        not available, leaving a clean correction point. \\
    \bottomrule
  \end{tabularx}
\end{table*}

\subsection{Contract objects and executable path}

A contract-grounded interaction is represented by the tuple
\[
\mathcal{I} =
\langle m,s,t,p,u,c,b,a,q,o\rangle .
\]
The symbols form three connected layers.  The interaction context consists of
the pinned model and hash \(m\) and scenario step \(s\).  The user-visible
referent and responsibility are the target \(t\) and provider \(p\), while
\(u\) and \(c\) record the context-specific capability use and its reusable
capability.  Execution and evidence are represented by runtime binding \(b\),
concrete action \(a\), assertions \(q\), and observations \(o\).  Thus, the
formal tuple answers the same questions exposed by the interface: what was
selected, who may respond, what may be executed, and what evidence justifies
presentation.

In the running example, \(t\) is \texttt{brand\_panel3}, \(p\) resolves to the
Reception Agent, \(s,u,c\) describe grounded answering, \(b,a\) map that
capability to backend chat, and \(q,o\) retain the target and provider evidence
required before the answer can be shown.

The required declaration path is
\[
s \rightarrow u \rightarrow c \rightarrow b \rightarrow a.
\]
A scenario step cannot invoke action \(a\) directly.  Its capability use \(u\)
must name the provider and targets and resolve exactly one capability, binding,
and action; a validation case covers the same binding and assertions.  This
prevents a generic \texttt{POST /chat} invocation from bypassing the declaration
that the Reception Agent may answer about this panel.  It also supplies a
traceable answer to ``why this operation?'': the operation is not selected by
endpoint type alone but by the target-specific scenario, provider, and
capability relation.

The provider \(p\) is resolved independently of display names.  For target
\(t\), the resolver selects the most specific non-empty responsibility tier:
\[
R(t) =
\begin{cases}
R_{\mathrm{asset}}(t), & \text{if non-empty},\\
R_{\mathrm{group}}(t), & \text{otherwise if non-empty},\\
R_{\mathrm{zone}}(t), & \text{otherwise}.
\end{cases}
\]
Exactly one agent at the selected tier is required; several owners are
ambiguous and no owner leaves the target unassigned.  The rule is a policy of
specificity, not a psychological model of whom visitors will intuitively expect
to answer: an explicit asset-level assignment overrides broader group and zone
defaults.  For \texttt{brand\_panel3}, this selects the Reception Agent.  If the
Lounge Host is active, one declared handoff may reach that provider; longer
routes remain diagnostic only.

This distinction has a direct interaction-design consequence.  A deterministic
resolver can still appear arbitrary when its provider differs from the active,
nearest, or most visually salient agent.  The interface should therefore expose
the selected provider and the reason for a handoff instead of asking users to
infer the asset--group--zone policy.  The contract supplies that explanation
state; the user study examines whether the current cues actually make the
provider transition and division of responsibility understandable.

\subsection{Admission and response conditions}

A deictic request is admitted only if the following conditions hold.  Each gate
is both a technical check and a possible explanation or recovery point:

\begin{enumerate}
  \item \textbf{Environment consistency---is this still the same scene?}
  The session is pinned to \(m\), and the supplied model hash equals the current
  project hash and stored contract fingerprint.

  \item \textbf{Reference clarity---what object was selected?}
  The supplied entity and source identifier resolve to the same unique scene
  asset in \(m\), with a resolved, non-ambiguous selection state.

  \item \textbf{Responsibility clarity---why may this agent answer?}
  The responsibility rule yields exactly one provider \(p\), and \(p\) is
  either the active agent or one permitted direct handoff away.

  \item \textbf{Capability clarity---what may the agent do here?}
  One capability use \(u\) for the trusted target and provider resolves to one
  capability \(c\), binding \(b\), and action \(a\).

  \item \textbf{Executable request---can this operation run safely?}
  The request conforms to the declared action schema.
\end{enumerate}

For the running example, these gates confirm the pinned room and panel identity,
resolve the Reception Agent, permit the Lounge Host's direct handoff, and select
the grounded-answer action before the response model is called.  The generated
proposal is then compared with the admitted path: handoffs must be declared and
all returned identifiers must agree.  A match is committed with evidence;
another target or undeclared agent rolls back provisional state.  Later
observations evaluate the assertions as pass, fail, inconclusive, or error.

Figure~\ref{fig:admission-boundary} shows the order and separates the technical
gates from their interface consequences.  The current cues use lightweight
visual status and short guidance text; the contract itself does not require
that particular presentation.

\begin{figure*}[t]
  \centering
  \begin{tikzpicture}[
    font=\footnotesize,
    stage/.style={
      draw,
      rounded corners=2pt,
      align=center,
      text width=27mm,
      minimum height=17mm,
      inner sep=1.5mm
    },
    gate/.style={
      stage,
      double,
      double distance=0.6pt
    },
    reject/.style={
      draw,
      rounded corners=2pt,
      align=center,
      text width=38mm,
      minimum height=14mm,
      inner sep=1.5mm,
      dashed
    },
    flow/.style={-{Latex[length=2.2mm]},semithick},
    fail/.style={-{Latex[length=2.2mm]},dashed},
    label/.style={font=\scriptsize,align=center,fill=white,inner sep=1pt}
  ]
    \node[stage] (select) at (0,0) {\textbf{1. Target selection}\\
      stable entity/source ID\\
      \textit{UI: spatial highlight}};
    \node[gate] (identity) at (3.35,0) {\textbf{Gate 1: identity}\\
      pinned hash, unique asset, no ambiguity};
    \node[gate] (route) at (6.70,0) {\textbf{Gate 2: responsibility}\\
      asset $>$ group $>$ zone; local or direct};
    \node[gate] (binding) at (10.05,0) {\textbf{Gate 3: capability}\\
      exact scenario, provider, binding, action};
    \node[stage] (model) at (13.40,0) {\textbf{Response model}\\
      called only after deterministic preflight};

    \draw[flow] (select) -- (identity);
    \draw[flow] (identity) -- (route);
    \draw[flow] (route) -- (binding);
    \draw[flow] (binding) -- (model);

    \node[reject] (reject) at (4.85,-3.25) {\textbf{Clarify or guide}\\
      expose target, responsibility, or capability problem; preserve a clean correction point};
    \draw[fail] (identity.south) to[bend right=13]
      node[label,left,pos=.52] {stale or ambiguous} (reject.north west);
    \draw[fail] (route.south) --
      node[label,right,pos=.48] {no unique/permitted provider} (reject.north);
    \draw[fail] (binding.south) to[bend left=13]
      node[label,right,pos=.52] {zero/multiple actions} (reject.north east);

    \node[gate] (evidence) at (13.40,-3.25) {\textbf{Gate 4: evidence}\\
      target, provider, action, and handoff agreement};
    \node[stage,text width=35mm] (present) at (9.30,-3.25)
      {\textbf{2. Commit and present}\\
      render grounded answer and responsible-agent state};
    \draw[flow] (model) -- (evidence);
    \draw[flow] (evidence) -- (present);

    \node[reject] (rollback) at (13.40,-5.65) {\textbf{Rollback and explain}\\
      discard mismatching proposal; retain pending target; report unavailable handoff};
    \draw[fail] (evidence) --
      node[label,left] {mismatch or undeclared handoff} (rollback);
  \end{tikzpicture}
  \caption{Fail-closed interaction boundary and its transparency points.  The
  backend re-resolves identity, responsibility, and capability before invoking
  the response model.  Each rejection can be translated into a specific
  clarification or guidance state instead of a generic error.  After
  generation, evidence agreement determines whether the answer is presented or
  rolled back.  The shown highlight, handoff, and guidance cues are the current
  prototype mapping, not requirements of the formal contract.}
  \Description{A two-row pipeline begins with target selection and a spatial
  highlight.  Identity, responsibility, and capability gates precede the
  response model and point to a clarification and guidance state when they
  fail.  The response then passes through an evidence gate before commit and
  presentation.  A mismatching response or undeclared handoff points to a
  rollback and explanation state.}
  \label{fig:admission-boundary}
\end{figure*}

\subsection{Guarantee boundary and interaction claims}

The implementation provides a \emph{model-consistency guarantee}, together
with structured evidence that an interface may use for transparency.  Its
boundary is deliberately narrower than semantic correctness or improved trust:

\begin{itemize}
  \item \textbf{System consistency.} If the example is admitted, one trusted
  panel resolves to one declared provider reachable by at most one handoff and
  one capability-to-action binding.  The response is presented only when its
  target, provider, handoff, binding, and action evidence agree with that path.

  \item \textbf{Transparency and control support.} The same relations identify
  what the system selected, why a provider was chosen, where a handoff occurred,
  and which gate rejected a request.  This supports specific feedback,
  explanation, and correction states instead of silent fallback.

  \item \textbf{Not guaranteed.} The contract does not establish that the
  visitor intended the selected collider, that the provider assignment matches
  every visitor's expectation, that generated content is factually correct, or
  that a particular visual cue improves comprehension, usability, or calibrated
  trust.
\end{itemize}

The interaction-design hypothesis is therefore not that adding a badge or an
explanation automatically creates trust.  It is that an explicit contract
provides the minimum state needed to communicate target, responsibility,
handoff, capability, and rejection without inventing a post-hoc rationale.  Our
user study tests this concrete mapping by examining target-selection clarity,
answer--target fit, handoff observation and understanding, final-agent clarity,
perceived responsibility, predictability, guidance, and overall ease.  This
connects the formal guarantee to an empirical HCI question: not only whether the
system follows its declared path, but whether users can understand the parts of
that path that the interface exposes.
